% !TeX program = lualatex
% !TeX root    = UCIMED-Evaluacion-plantilla.tex
%%%%%%%%%%%%%%%%%%%%%%%%%%%%%%%%%%%%%%%%%%%%%%%%%%%%%%%%%%%%%%%%%%%%%%%%%%%%%%%
% PLANTILLA de evaluación UCIMED — documento para trabajo colaborativo.
%
% Archivos que deben acompañar a este .tex en la misma carpeta:
%     format-UCIMED.sty   colors-UCIMED.sty   comands-UCIMED.sty
%     eval-UCIMED.sty     headers-UCIMED.sty
% (el logo institucional va incrustado en comands-UCIMED.sty; no hace falta PNG)
%
% Compilar con lualatex, DOS pasadas (o `latexmk -lualatex`), porque el puntaje
% total del encabezado y los subtotales por sección viajan por el .aux.
%
%   solutions = true   → solucionario: imprime respuestas y marca la correcta.
%   solutions = false  → versión de aplicación para el estudiantado.
%%%%%%%%%%%%%%%%%%%%%%%%%%%%%%%%%%%%%%%%%%%%%%%%%%%%%%%%%%%%%%%%%%%%%%%%%%%%%%%
\documentclass[11pt,twoside]{article}

  \usepackage{format-UCIMED}     % preámbulo mínimo (lualatex)
  \usepackage{colors-UCIMED}     % paleta institucional
  \usepackage{comands-UCIMED}    % \setclass \setcycle \setgroup + logo

%==< Datos del documento >=====================================================
  \title{Prueba Parcial}
  \author{Luis Fernando Umaña Castro}
  \date{13 de agosto de 2026}

  \setclass{CB-09}      % CB-09 | CBI-06 | CBI-07  (fija curso e institución)
  \setcycle{2026-II}
  \setgroup{01}         % en documentos multiversión, \setgroup = versión

  \usepackage{eval-UCIMED}       % \question \qpart \pts \consolidatePoints …
  \setpartialnumber{1}           % 1 → «Prueba Parcial I»;  0 → sin numeral
  \usepackage{headers-UCIMED}    % \pagestyle{partial} y {uplain}
  \pagestyle{uplain}

%==< Banderas de compilación >=================================================
  \setbool{test}{true}           % formato de examen (ítems numerados)
  \setbool{solutions}{true}      % ← cambiar a false para la versión del examen

%%%%%%%%%%%%%%%%%%%%%%%%%%%%%%%%%%%%%%%%%%%%%%%%%%%%%%%%%%%%%%%%%%%%%%%%%%%%%%%
\begin{document}
%%%%%%%%%%%%%%%%%%%%%%%%%%%%%%%%%%%%%%%%%%%%%%%%%%%%%%%%%%%%%%%%%%%%%%%%%%%%%%%
\thispagestyle{partial}

\ifthenelse{\boolean{solutions}}{%
  \begin{center}\large\textbf{Solucionario}\hfill\thegroup\end{center}
}{%
  \AANlines
  \noindent\makebox[\textwidth]{\rule{10cm}{0.4pt}}

  \section*{Instrucciones generales}
  \begin{enumerate}
    \item Utilice bolígrafo de color negro o azul. No se aceptan respuestas en
          lápiz y se calificarán con cero.
    \item No utilice corrector. Si lo hace se calificará con cero.
    \item Dispone de 100 minutos para resolver esta prueba y sólo se atenderán
          consultas generales durante los primeros 30 minutos.
    \item Durante esta prueba no se permite el uso de aparatos de
          radiocomunicación.
  \end{enumerate}
  \newpage
}

%%%%%%%%%%%%%%%%%%%%%%%%%%%%%%%%%%%%%%%%%%%%%%%%%%%%%%%%%%%%%%%%%%%%%%%%%%%%%%%
%==< Sección 1 · Complete >====================================================
%%%%%%%%%%%%%%%%%%%%%%%%%%%%%%%%%%%%%%%%%%%%%%%%%%%%%%%%%%%%%%%%%%%%%%%%%%%%%%%
\section*{Complete}
En esta parte se incluyen 2 preguntas de complete.
\textbf{Puntaje Total = \ref*{CT1} puntos.}

\vspace{5mm}
\noindent\textbf{Instrucciones}: escriba en el espacio en blanco la palabra que
completa correctamente cada frase. Dispone únicamente del grupo de conceptos
dado; \textbf{no se repiten conceptos}.

\vspace{5mm}
\begin{enumerate}[1.]
  \item La presión que ejerce una columna de fluido en reposo depende de su
        densidad, de la gravedad y de la \completeline{profundidad} \upt[add]
  \item El cociente entre el esfuerzo y la deformación en el régimen elástico
        se llama módulo de \completeline{Young} \upt[add]
\end{enumerate}
\consolidatePoints{CT1}

\newpage
%%%%%%%%%%%%%%%%%%%%%%%%%%%%%%%%%%%%%%%%%%%%%%%%%%%%%%%%%%%%%%%%%%%%%%%%%%%%%%%
%==< Sección 2 · Selección única >=============================================
%%%%%%%%%%%%%%%%%%%%%%%%%%%%%%%%%%%%%%%%%%%%%%%%%%%%%%%%%%%%%%%%%%%%%%%%%%%%%%%
\section*{Selección única}
En esta parte se incluye 1 pregunta de selección única.
\textbf{Puntaje Total = \ref*{SU1} puntos.}

\vspace{5mm}
\noindent\textbf{Instrucciones}: marque con una ``\textbf{X}'' la letra que
acompaña la respuesta correcta.

\begin{enumerate}[1.]
% --- Un ítem completo. En producción cada ítem vive en su propio archivo y se
%     trae con \input{...}; aquí va en línea sólo para mostrar la estructura.
\question{
  \pts[add]{2}
  Un recipiente contiene agua (\(\rho = \SI{1.00e3}{\kilo\gram\per\cubic\meter}\)).
  ¿Cuál es la presión manométrica a una profundidad de \SI{1.50}{\meter}?
  \begin{enumerate}[a)]
    \item \SI{1.47}{\kilo\pascal}
    \item \SI{14.7}{\kilo\pascal} \rightoption
    \item \SI{147}{\kilo\pascal}
    \item \SI{1.01e5}{\pascal}
    \item \SI{9.81}{\kilo\pascal}
  \end{enumerate}
}{}
\end{enumerate}
\consolidatePoints{SU1}

\newpage
%%%%%%%%%%%%%%%%%%%%%%%%%%%%%%%%%%%%%%%%%%%%%%%%%%%%%%%%%%%%%%%%%%%%%%%%%%%%%%%
%==< Sección 3 · Desarrollo >==================================================
%%%%%%%%%%%%%%%%%%%%%%%%%%%%%%%%%%%%%%%%%%%%%%%%%%%%%%%%%%%%%%%%%%%%%%%%%%%%%%%
\section*{Desarrollo}
A continuación se incluye 1 problema de desarrollo.
\textbf{Puntaje Total = \ref*{DE1} puntos.}

\vspace{5mm}
\noindent\textbf{Instrucciones}: resuelva de manera clara y ordenada, mostrando
todos los pasos.

\partialenumi
\begin{enumerate}
  \setthisindent
% --- Ítem de desarrollo con sub-partes.
%     Las sub-partes van en el PRIMER argumento de \question, dentro de un
%     enumerate[a)]; el segundo argumento queda vacío.
\question{
  Una arteria de radio \(r = \SI{2.00}{\milli\meter}\) transporta sangre con un
  caudal volumétrico \(Q = \SI{5.00}{\milli\liter\per\second}\).
  \begin{enumerate}[a)]
    \qpart{Calcule la rapidez media del flujo. \pts[add]{3}}{
      De la ecuación de continuidad, \(Q = A v\), con
      \(A = \pi r^2 = \SI{1.26e-5}{\square\meter}\):
      \[ v = \frac{Q}{A}
           = \frac{\SI{5.00e-6}{\cubic\meter\per\second}}{\SI{1.26e-5}{\square\meter}}
           = \SI{0.398}{\meter\per\second} \]
    }
    \qpart{Si una placa reduce el radio a la mitad, ¿en qué factor cambia la
           rapidez? \pts[add]{2}}{
      Como \(A \propto r^2\) y \(Q\) se conserva, \(v \propto 1/r^2\): al
      reducirse el radio a la mitad, la rapidez se cuadruplica
      (\(v' = 4v = \SI{1.59}{\meter\per\second}\)).
    }
  \end{enumerate}
}{}
\end{enumerate}
\consolidatePoints{DE1}

%%%%%%%%%%%%%%%%%%%%%%%%%%%%%%%%%%%%%%%%%%%%%%%%%%%%%%%%%%%%%%%%%%%%%%%%%%%%%%%
\end{document}
%%%%%%%%%%%%%%%%%%%%%%%%%%%%%%%%%%%%%%%%%%%%%%%%%%%%%%%%%%%%%%%%%%%%%%%%%%%%%%%
